\section{Metodología de desarrollo}
%  TODO - re Re RE RE dactar
El desarrollo de este trabajo se ha llevado a cabo siguiendo una metodología basada en la construcción iterativa de prototipos, tal y como se estableció en la propuesta inicial del trabajo (TFT01). Esta elección metodológica se justifica por la naturaleza exploratoria del proyecto: al abordar la implementación directa de protocolos de comunicación vehicular sobre hardware de propósito general, sin disponer de precedentes directos en el grupo de trabajo, resultaba imprescindible disponer de un enfoque que permitiera validar hipótesis técnicas de forma temprana y ajustar el diseño en función de los resultados obtenidos en cada iteración.

La metodología basada en prototipos se estructura en cuatro fases principales: análisis de requisitos, diseño rápido y construcción del prototipo inicial, evaluación y refinamiento iterativo, y documentación y presentación. Estas fases no son estrictamente secuenciales sino que admiten solapamiento parcial: en particular, las fases de construcción y evaluación se ejecutaron de forma alternada durante la mayor parte del desarrollo, con ciclos cortos de implementación seguidos de pruebas sobre el banco de ensayo Arduino.

Este enfoque permitió, por ejemplo, detectar de forma temprana las limitaciones de la primera implementación del protocolo ISO-TP, que no contemplaba la gestión correcta del \textit{Flow Control} en mensajes multi-frame, y corregirlas antes de proceder al desarrollo de las capas superiores. De igual modo, la decisión de incorporar el módulo de simulación de ruido CAN en el simulador Arduino surgió como respuesta a los resultados de las primeras pruebas de integración, en las que se constató la necesidad de validar la robustez del sistema frente a condiciones de bus no ideales.


% TODO - revisar y redactar
\section{Fases y tareas}

\subsection{Fase 1 — Análisis de requisitos \textnormal{(50 horas estimadas)}}

Esta fase inicial se dedicó al estudio de los fundamentos técnicos necesarios para abordar el desarrollo del sistema.

\textbf{Tarea 1.1 — Estudio de la plataforma hardware.} Se analizaron las características de la Raspberry Pi como plataforma de ejecución de la herramienta de diagnóstico, incluyendo el soporte de la interfaz SocketCAN en Linux, la compatibilidad con el controlador CAN MCP2515 y las opciones disponibles para el servidor BLE. Asimismo, se evaluó el Arduino MKR como plataforma del simulador, estudiando la librería CAN disponible y sus capacidades de temporización.

\textbf{Tarea 1.2 — Estudio de los protocolos de comunicación CAN.} Se realizó un análisis detallado de los estándares ISO 11898 (CAN), ISO 15765-2 (ISO-TP) y SAE J1979 (OBD-II), prestando especial atención a la estructura de las tramas, los mecanismos de control de flujo y el formato de codificación de cada parámetro. Se revisaron asimismo los conceptos fundamentales de UDS (ISO 14229-1) relevantes para la implementación de las respuestas negativas.

\textbf{Tarea 1.3 — Estudio de Bluetooth Low Energy y librerías disponibles.} Se estudiaron los conceptos fundamentales de BLE GATT, el servicio Nordic UART Service (NUS) y las opciones de implementación disponibles en Python para la Raspberry Pi, evaluando las librerías \texttt{bleak} y \texttt{bless}. En el lado de la aplicación móvil, se analizó la librería \texttt{react-native-ble-plx} y su compatibilidad con las características de escritura y notificación del servicio NUS.

\textbf{Tarea 1.4 — Análisis de herramientas software.} Se definió el conjunto de herramientas de desarrollo a utilizar: Python 3 con las librerías \texttt{python-can} y \texttt{can-isotp} para la capa de transporte CAN, SQLite para la persistencia de datos, React Native con Expo para la aplicación móvil, y el IDE de Arduino para el simulador. Se estableció asimismo la estructura del repositorio y el flujo de trabajo con control de versiones Git.

\subsection{Fase 2 — Diseño rápido y construcción \textnormal{(75 horas estimadas)}}

\textbf{Tarea 2.1 — Diseño de la arquitectura del sistema.} Se definió la arquitectura general del sistema en tres capas: simulador Arduino, herramienta de diagnóstico en Python y aplicación móvil. Para la herramienta Python se diseñó una arquitectura modular basada en interfaces abstractas (\texttt{ITransport}, \texttt{IProtocolBuilder}, \texttt{IDataDecoder}, \texttt{IDiagnosticSession}), con inyección de dependencias para facilitar las pruebas. Se definieron asimismo el esquema de la base de datos SQLite y el protocolo de comandos JSON utilizado sobre BLE.

\textbf{Tarea 2.2 — Implementación del primer prototipo funcional.} Se desarrolló una primera versión operativa del sistema que incluía: el simulador Arduino respondiendo a peticiones OBD-II básicas en modo 0x01 (RPM, temperatura de refrigerante, velocidad), la herramienta Python con transporte ISO-TP funcional y sesión de diagnóstico básica, y una primera versión de la aplicación móvil con pantalla de conexión BLE y visualización de datos en tiempo real.

\subsection{Fase 3 — Evaluación, refinamiento e implementación \textnormal{(125 horas estimadas)}}

Esta fue la fase de mayor duración e intensidad técnica, en la que el prototipo inicial evolucionó de forma iterativa hasta el sistema final.

\textbf{Tarea 3.1 — Extensión del simulador Arduino.} Se amplió el simulador con soporte completo para los cuatro modos OBD-II (0x01, 0x03, 0x04, 0x09), implementando la respuesta VIN mediante la secuencia ISO-TP multi-frame (FF + CF). Se añadió el modelo físico de simulación del motor, con la lógica de aceleración, velocidad, temperaturas y consumo de combustible. Posteriormente se incorporaron los mecanismos de simulación de ruido CAN y la interrupción hardware para inyección aleatoria de DTCs.

\textbf{Tarea 3.2 — Extensión de la herramienta Python.} Se completó la herramienta con el módulo de logging en SQLite, el decorador \texttt{LoggingTransport}, el monitor de datos en tiempo real con hilo dedicado, la sesión de diagnóstico extendida (\texttt{LoggedDiagnosticSession}) y el servidor BLE GATT con el manejador de comandos JSON (\texttt{BtCommandHandler}).

\textbf{Tarea 3.3 — Desarrollo completo de la aplicación móvil.} Se implementaron las cinco pantallas de la aplicación (Connection, Dashboard, DTCs, Logs, Console), los almacenes de estado con Zustand, el adaptador BLE con reintentos y gestión de timeout, y el adaptador simulado (\texttt{MockAdapter}) para el desarrollo sin hardware.

\textbf{Tarea 3.4 — Pruebas de integración.} Se realizaron pruebas del sistema completo sobre el banco de ensayo, verificando el comportamiento ante los escenarios de ruido CAN configurados en el simulador, la correcta persistencia de los datos en SQLite y la estabilidad de la conexión BLE durante sesiones de monitorización prolongadas.

\subsection{Fase 4 — Documentación y presentación \textnormal{(50 horas estimadas)}}

\textbf{Tarea 4.1 — Redacción de la memoria.} Redacción del presente documento, incluyendo la elaboración de las tablas y figuras de los protocolos de comunicación, la documentación del código fuente incluida en los anexos y la preparación del material de apoyo para la defensa.


% TODO - revisar y redactar
\section{Desviaciones respecto a la planificación inicial}

La planificación inicial recogida en el TFT01 contemplaba una distribución de 300 horas totales, con especial peso en la fase de evaluación e implementación (125 horas). En líneas generales, esta estimación resultó ajustada a la realidad del desarrollo, si bien se produjeron algunas desviaciones significativas que conviene señalar.

La implementación del protocolo ISO-TP requirió un esfuerzo mayor al previsto inicialmente, en particular en lo referente a la gestión del \textit{Flow Control} y a la sincronización entre el envío de la secuencia FF+CF en el simulador Arduino y la recepción en la Raspberry Pi. Esta fase consumió aproximadamente un 20\% más de tiempo del estimado, que se compensó reduciendo el alcance inicial de la aplicación móvil, que en la propuesta contemplaba también visualización gráfica de tendencias temporales.

Por otra parte, la incorporación del módulo de simulación de ruido CAN no estaba prevista en la planificación original, pero se identificó como necesaria durante las primeras pruebas de integración. Su desarrollo se realizó dentro de la Fase 3, absorbiendo parte del tiempo reservado a las pruebas sobre el vehículo real, que finalmente no pudieron llevarse a cabo por las razones expuestas en la Sección~\ref{sec:alcance}.

La decisión de estructurar el código Python con una arquitectura de interfaces y capas desacopladas supuso una inversión inicial de tiempo superior a la contemplada en el diseño rápido, pero resultó rentable en las fases posteriores al facilitar considerablemente la extensión del sistema y la realización de pruebas con el transporte simulado.
